\scanpage{54}
\mspage{1}

\uline{\unsure{11 1.84|11 I.84}}

\uline{Chirurgie conforme}

\textcircled{\unsure{4|A|1}} Surface conforme à bord. On \deleted{prend des} la définit par des cartes locales qui, au voisinage d'un point du bord, sont modelées sur \inserted{$\overline H=$} $(\Im z\geq 0)\subset\mathbb C$. Dans le cas $X$ orienté i.e structure holomorphe, la structure peut \deleted{\unsure{se|de}} définir par un faisceau de $\mathbb C$-algèbres, qui, sur le modèle $\overline H$ précédent, aux points $x$ de $\partial\overline H=\mathbb R$, est formée de germes de fonctions qui se \uline{prolongent} \deleted{\unsure{alors|des}} fonctions hol. \deleted{aux pts de} \inserted{au voisinage} de $x$ \uline{dans} $C$.

À ne pas confondre avec le sous faisceau plus grand du faisceau des fonctions complexes continues, savoir \deleted{celui des} \inserted{le f. des} fonctions continues qui sont hol dans $\operatorname{Int}(X)$.

\uline{Théorème 1} La structure conforme de $X$ (surface conforme à bord) est connue quand on la connaît dans $\operatorname{Int}(X)$, (ce qui suffit à déterminer le faisceau des fonctions \inserted{complexes} continues qui sont conformes dans $\operatorname{Int}(X)$).

Il suffit de voir que \unsure{cel|cela} permet de réconstituer les \uline{cartes} conformes des voisinages des points du bord. Je dis que ces cartes
% archive sheet number: 56
